\section{Технический проект}
\subsection{Архитектура онлайн-площадки для сдачи вещей в аренду}

Необходимо спроектировать и разработать веб-приложение для онлайн площадки для сдачи вещей в аренду. 
Проектируемая программа представляет собой веб-приложение, построенное по архитектуре «клиент – сервер». Клиентская часть представлена веб-интерфейсом на HTML/CSS/JavaScript, который открывается в браузере пользователя. Серверная часть реализована на языке Python с использованием веб-фреймворка Django и предоставляет REST API для чтения и изменения данных. Хранение структурированных данных выполняется в СУБД PostgreSQL, кэширование и управление WebSocket-соединениями — в Redis. Развёртывание всех компонентов серверной части (Django, PostgreSQL, Redis, Nginx) выполняется с использованием системы контейнеризации Docker, что обеспечивает идентичность окружений на этапах разработки, тестирования и эксплуатации. 

Общая схема взаимодействия компонентов приведена на рисунке \ref{uml:image}

\begin{figure}[h]
\center{\includegraphics[width=0.7\linewidth]{canvas}}
\caption{Общая схема взаимодействия компонентов}
\label{uml:image}
\end{figure}

Выбор такой архитектуры обоснован требованиями проекта: необходимо обеспечить одновременную работу пользовательской части (просмотр товаров, бронирование, чаты) и административной части с единой базой данных. Сервер централизует бизнес-логику, валидацию, контроль пересечения дат бронирования, а также статусную модель аренды.

\subsection{Обоснование выбора технологий проектирования}

Для реализации клиентской части выбраны HTML5, CSS3 и Vanilla JavaScript. Использование стандартных веб-технологий позволяет создать полноценный интерфейс без подключения крупных frontend-фреймворков. Для данного проекта такой подход является достаточным, поскольку приложение имеет понятную структуру: каталог вещей, формы, личный кабинет, бронирования, чаты и уведомления.

HTML5 используется для описания структуры интерфейса. CSS3 применяется для оформления страниц, адаптивной верстки и визуального выделения интерактивных элементов. JavaScript отвечает за динамическое изменение интерфейса, обработку действий пользователя, отправку запросов к серверу и работу с WebSocket.

Для взаимодействия с сервером используется Fetch API. Он встроен в современные браузеры и позволяет отправлять HTTP-запросы без дополнительных библиотек. В проекте все запросы централизованы через модуль api.js, что позволяет единообразно обрабатывать ответы сервера, ошибки и авторизацию.

Для хранения данных сессии применяется localStorage. В нем сохраняются JWT-токены, идентификатор пользователя, имя, email и тип учетной записи. Это позволяет сохранять состояние авторизации после перезагрузки страницы.

Для выбора периода аренды используется библиотека flatpickr. Она подключается к интерфейсу выбора дат и обеспечивает удобный ввод даты начала и окончания аренды. Использование готовой библиотеки снижает риск ошибок при ручном вводе дат и улучшает пользовательский опыт.

\subsection{Архитектура клиентской части}

Клиентская часть веб-приложения реализована как одностраничное приложение. Пользователь работает с одним HTML-документом, а смена экранов выполняется средствами JavaScript без полной перезагрузки страницы. Такой подход позволяет сделать интерфейс более быстрым и удобным, так как при переходе между разделами обновляется только содержимое рабочей области приложения.

Клиентская часть состоит из основного файла index.html, общего файла стилей styles.css и набора JavaScript-модулей, расположенных в папке js. Такое разделение позволяет отделить структуру страницы, внешний вид интерфейса и программную логику приложения. Ниже приведена схема компонентов клиентской части \ref{clientgramm:image}

\begin{figure}[H]
\center{\includegraphics[width=1\linewidth]{clientgramm}}
\caption{Схема компонентов клиентской части}
\label{clientgramm:image}
\end{figure}

\subsubsection{Файл index.html}
Файл index.html выступает точкой входа в приложение. В нем размещается базовая HTML-структура страницы, подключаются таблица стилей, сторонние библиотеки и JavaScript-файлы. Также в нем находится основной контейнер, в который с помощью JavaScript подставляется содержимое текущего раздела сайта.

\subsubsection{Файл style.css}
Файл styles.css содержит все стили клиентской части. В нем описывается внешний вид шапки сайта, карточек вещей, кнопок, форм, модальных окон, страниц профиля, бронирований, чатов и уведомлений. Кроме того, в этом файле задаются правила адаптивной верстки, благодаря которым интерфейс корректно отображается на компьютерах, планшетах и мобильных устройствах.

\subsubsection{Файл config.js}
JavaScript-код разделен на несколько файлов по функциональному назначению. Файл config.js хранит общие настройки приложения, например базовый адрес API и адрес WebSocket-соединений. Файл api.js отвечает за отправку запросов к серверу, добавление токена авторизации в заголовки, обработку ошибок и обновление access-токена.

\subsubsection{Файл auth.js}
Файл auth.js содержит логику регистрации, входа, выхода из аккаунта, восстановления пароля и проверки состояния авторизации. В нем также выполняется сохранение и удаление данных пользователя в localStorage.

\subsubsection{Файл catalog.js}
Файл catalog.js отвечает за работу каталога 
вещей. В нем реализована загрузка списка товаров, отображение карточек, получение категорий, фильтрация, пагинация и открытие подробной информации о выбранной вещи.

\subsubsection{Файл bookings.js}
Файл bookings.js содержит функции, связанные с арендой и бронированиями. Он отвечает за выбор периода аренды, создание заявки, получение входящих и исходящих бронирований, а также подтверждение или отмену заявок.

\subsubsection{Файл chats.js}
Файл chats.js реализует работу с чатами. В нем находится логика получения списка диалогов, открытия переписки, отправки сообщений и изображений, а также подключения к WebSocket для получения новых сообщений в реальном времени.

\subsubsection{Файл profile.js}
Файл profile.js отвечает за личный кабинет пользователя. Через него выполняется загрузка данных профиля, отображение информации о пользователе, редактирование профиля и изменение пароля.

\subsubsection{Файл notifications.js}
Файл notifications.js содержит функции для работы с уведомлениями. Он отвечает за получение уведомлений с сервера, отображение счетчиков, обработку прочитанных уведомлений и подключение к WebSocket-каналу уведомлений.

\subsubsection{Файл contracts.js}
Файл contracts.js отвечает за работу с договорами аренды и демонстрационным процессом оплаты. В нем реализовано получение договора по бронированию, открытие договора, подписание и дальнейшее выполнение демонстрационной оплаты.

\subsubsection{Файл main.js}
Файл main.js является основным модулем запуска приложения. Он инициализирует интерфейс, подключает обработчики событий, проверяет состояние авторизации пользователя, загружает начальные данные и отображает главную страницу сайта.

\subsubsection{Связь с API}
Связь клиентской части с сервером выполняется через Django REST API. Для обычных операций используются HTTP-запросы к endpoint с базовым адресом api, а для работы в реальном времени применяются WebSocket-соединения. Такой подход позволяет совместить классическое API-взаимодействие с мгновенным получением сообщений и уведомлений.

\subsection{Реализация ключевых функциональных модулей}
\subsubsection{Модуль каталога вещей}

Модуль каталога реализован в файле catalog.js. Он отвечает за загрузку и отображение вещей, доступных для аренды. При открытии главной страницы модуль отправляет запрос к endpoint /api/items/, получает список объектов и формирует карточки вещей.

Каждая карточка содержит название, краткое описание, изображение, цену аренды за день, рейтинг и кнопки действий. Пользователь может открыть подробную информацию о вещи, перейти к аренде или начать чат с владельцем.

Также в модуле каталога реализована работа с категориями. Категории загружаются с сервера через API, после чего пользователь может выбрать нужную категорию. При выборе категории список вещей обновляется.

Дополнительно реализована фильтрация по доступности на выбранный период. Пользователь указывает даты аренды, после чего frontend передает параметры фильтрации серверу и отображает только подходящие вещи.

\subsubsection{Модуль карточки вещи}

Карточка вещи также относится к модулю catalog.js, так как является частью работы с каталогом. При выборе товара открывается модальное окно с подробной информацией. В нем отображаются изображения, описание, владелец, стоимость аренды, рейтинг и отзывы.

Из модального окна пользователь может начать процесс аренды или открыть чат с владельцем. Если текущий пользователь является владельцем вещи, интерфейс отображает дополнительные действия: редактирование или удаление объявления.

Такой подход позволяет пользователю получить всю основную информацию о вещи без перехода на отдельную страницу и без перезагрузки приложения.

\subsubsection{Модуль аренды и бронирования}

Модуль аренды и бронирования реализован в файле bookings.js. Он отвечает за выбор периода аренды, создание заявки и управление статусами бронирований.

При открытии окна аренды пользователь выбирает дату начала и дату окончания аренды. Для выбора дат используется библиотека flatpickr. После выбора периода frontend рассчитывает параметры аренды и отправляет запрос на создание бронирования через endpoint /api/bookings/.

В модуле также реализовано отображение входящих и исходящих бронирований. Входящие бронирования показывают заявки на вещи, принадлежащие текущему пользователю. Исходящие бронирования показывают заявки, созданные самим пользователем.

Владелец вещи может подтвердить или отменить заявку. Пользователь также может отслеживать статус бронирования: ожидает подтверждения, подтверждено, отменено или завершено.

\subsubsection{Модуль авторизации и регистрации}

Модуль авторизации расположен в файле auth.js. Он отвечает за регистрацию, вход, выход, восстановление пароля и хранение данных сессии.

При входе пользователь отправляет логин и пароль на endpoint /api/token/. Сервер возвращает access и refresh токены. Эти токены сохраняются в localStorage и используются при последующих защищенных запросах.

Регистрация выполняется через endpoint /api/users/register/. При регистрации пользователь указывает основные данные учетной записи и тип пользователя. В системе поддерживаются физическое лицо, индивидуальный предприниматель и юридическое лицо.

При выходе из аккаунта модуль очищает данные пользователя и токены из localStorage, после чего интерфейс возвращается в состояние неавторизованного пользователя.

\subsubsection{Модуль личного кабинета}

Модуль личного кабинета реализован в файле profile.js. Он отвечает за получение, отображение и изменение данных пользователя.

При открытии профиля frontend отправляет запрос к /api/users/profile/ и получает данные текущего пользователя. В интерфейсе отображаются имя пользователя, email, тип учетной записи и дополнительные поля.

Набор дополнительных полей зависит от типа пользователя. Для физического лица используются паспортные данные и ИНН. Для индивидуального предпринимателя или юридического лица могут использоваться ОГРНИП, КПП, название организации и другие реквизиты.

Также в модуле профиля реализована смена пароля. Пользователь вводит старый и новый пароль, после чего frontend отправляет соответствующий запрос на сервер.

\subsubsection{Модуль управления объявлениями}

Функции управления объявлениями связаны с модулями catalog.js и profile.js. Пользователь может создать новое объявление, отредактировать существующее или удалить свою вещь.

При создании объявления пользователь указывает название, описание, цену за день, категорию и изображения. Основные данные вещи отправляются на сервер через /api/items/, а изображения загружаются отдельно через endpoint /api/items/upload-image/.

При редактировании объявления frontend сначала получает текущие данные вещи, затем отображает их в форме. После изменения данных пользователь сохраняет результат, и обновленная информация отправляется на сервер.

Удаление объявления доступно только владельцу вещи. Перед удалением интерфейс отображает подтверждение, чтобы избежать случайного удаления.

\subsubsection{Модуль чатов}

Модуль чатов расположен в файле chats.js. Он обеспечивает переписку между пользователями сервиса. Пользователь может открыть список диалогов, перейти к конкретному чату и отправлять сообщения.

Для получения списка чатов используется endpoint /api/chats/conversations/. Для отправки сообщения используется \texttt{/api/chats/conversations/{id}/messages/}. Сообщения могут содержать как текст, так и изображение.

Для получения новых сообщений без перезагрузки страницы используется WebSocket-соединение по адресу /ws/chat/{chat\_id}/. Когда сервер отправляет новое сообщение, frontend добавляет его в окно переписки и обновляет состояние чата.

\subsubsection{Модуль уведомлений}

Модуль уведомлений реализован в файле notifications.js. Он отвечает за получение уведомлений, отображение счетчиков и обработку событий в реальном времени.

Уведомления связаны с основными действиями в системе: созданием бронирования, подтверждением или отменой заявки, оплатой, новым сообщением и новым отзывом. Для получения списка уведомлений используется endpoint /api/notifications/.

Для мгновенного получения новых уведомлений используется WebSocket-канал /ws/notifications/. При получении уведомления frontend может показать всплывающее сообщение, обновить счетчик в навигации и при необходимости перенаправить пользователя к нужному разделу.

\subsubsection{Модуль договоров и оплаты}

Модуль договоров реализован в файле contracts.js. Он связан с процессом подтвержденной аренды. После подтверждения бронирования пользователь может открыть договор, просмотреть его данные и выполнить демонстрационное подписание.

Договор привязан к бронированию и содержит информацию о вещи, владельце, арендаторе, сроке аренды и стоимости. В интерфейсе предусмотрена возможность открыть договор и выполнить подписание от имени участника аренды.

Также в модуле реализована демонстрационная оплата через СБП. После подписания договора создается платежная сессия, пользователю отображаются данные оплаты, после чего он может подтвердить демонстрационный платеж. После успешной оплаты статус бронирования обновляется.

\subsection{Управление состоянием на клиенте}

Управление состоянием в проекте реализовано без отдельной библиотеки. Состояние приложения хранится средствами JavaScript и браузера.

Данные авторизации сохраняются в localStorage. К ним относятся access-токен, refresh-токен, идентификатор пользователя, имя пользователя, email и тип учетной записи. Это позволяет приложению восстанавливать состояние пользователя после перезагрузки страницы.

Локальное состояние интерфейса хранится в переменных соответствующих модулей. Например, модуль каталога хранит текущую категорию и выбранные даты фильтрации, модуль чатов - активный диалог и WebSocket-соединение, модуль бронирований - выбранное бронирование и период аренды.

Обработка ошибок централизована в api.js. Если сервер возвращает ошибку авторизации, модуль пытается обновить access-токен через refresh-токен. Если обновить токен невозможно, пользователь выходит из системы. Остальные ошибки отображаются в интерфейсе через сообщения, пустые состояния или toast-уведомления.

\subsection{Пользовательский интерфейс и UX}

Интерфейс приложения выполнен в едином визуальном стиле. В styles.css описаны общие цвета, оформление карточек, кнопок, форм, модальных окон и навигации. За счет единого файла стилей все разделы приложения выглядят согласованно.

Главная страница ориентирована на быстрый просмотр доступных вещей. Пользователь видит категории, фильтры и карточки товаров. Карточка содержит ключевую информацию: изображение, название, цену, рейтинг и основные действия.

Формы авторизации, регистрации, профиля и создания объявления имеют единый стиль. Это упрощает взаимодействие пользователя с системой и делает интерфейс предсказуемым.

Модальные окна используются для действий, которые не требуют перехода на отдельную страницу: просмотр вещи, выбор периода аренды, подтверждение удаления, просмотр договора и демонстрационная оплата.

В интерфейсе используются состояния загрузки, сообщения об ошибках, пустые состояния, счетчики уведомлений и всплывающие toast-сообщения. Это помогает пользователю понимать, что происходит после каждого действия.

Адаптивная верстка позволяет использовать сайт на разных устройствах. Элементы интерфейса перестраиваются в зависимости от ширины экрана, карточки переносятся на новые строки, а формы и модальные окна сохраняют удобный размер.

\subsection{Взаимодействие с backend API}

Frontend взаимодействует с backend-частью через REST API. Базовый путь API задается в config.js и равен /api. Все запросы выполняются через функции модуля api.js, что обеспечивает единообразную обработку заголовков, ошибок и авторизации.

Основные endpoint, используемые frontend-частью:

GET /api/items/ - получение списка вещей.

POST /api/items/ - создание объявления.

GET /api/items/{id}/ - получение подробной информации о вещи.

POST /api/items/upload-image/ - загрузка изображения вещи.

DELETE /api/items/images/{id}/ - удаление изображения.

GET /api/items/categories/ - получение категорий.

POST /api/bookings/ - создание бронирования.

GET /api/bookings/ - получение списка бронирований.

POST /api/bookings/{id}/confirm/ - подтверждение бронирования.

POST /api/bookings/{id}/cancel/ - отмена бронирования.

POST /api/token/ - авторизация пользователя.

POST /api/token/refresh/ - обновление access-токена.

POST /api/users/register/ - регистрация пользователя.

GET /api/users/profile/ - получение профиля.

PUT /api/users/profile/ - обновление профиля.

GET /api/chats/conversations/ - получение списка диалогов.

POST /api/chats/conversations/ - создание диалога.

POST /api/chats/conversations/{id}/messages/ - отправка сообщения.

GET /api/notifications/ - получение уведомлений.

GET /api/contracts/ - получение договоров.

Для обмена событиями в реальном времени используются WebSocket-соединения. Чат подключается к адресу \texttt{/ws/chat/\{chat\_id\}/}, а уведомления - к \texttt{/ws/notifications/}. Благодаря этому новые сообщения и уведомления появляются в интерфейсе без ручного обновления страницы.

\subsection{Требования к запуску клиентской части}

Клиентская часть не требует отдельной сборки через Vite, Webpack или npm. Приложение состоит из HTML-файла, CSS-файла и набора JavaScript-модулей, которые подключаются в браузере.

Для корректной работы frontend-части необходимо, чтобы backend-приложение Django было запущено и предоставляло доступ к REST API, WebSocket-маршрутам и медиафайлам. Также необходимо, чтобы сервер корректно отдавал статические файлы клиентской части.

Для работы приложения требуется современный браузер с поддержкой HTML5, CSS3, JavaScript ES6+, Fetch API, LocalStorage и WebSocket. К таким браузерам относятся актуальные версии Google Chrome, Mozilla Firefox, Microsoft Edge и других современных браузеров.

Таким образом, клиентская часть проекта представляет собой модульное одностраничное приложение на HTML, CSS и JavaScript. Разделение кода на отдельные файлы повышает читаемость и сопровождаемость проекта, а взаимодействие с backend через REST API и WebSocket обеспечивает выполнение основных пользовательских сценариев сервиса аренды вещей: просмотр каталога, регистрацию и авторизацию, создание объявлений, бронирование, общение между пользователями, получение уведомлений, работу с договорами и демонстрационную оплату.